\section{Functions, parameters and predicates}
\label{sec:functions}

A function can be equivariant, depend on finitely many atoms, or have no
finite support. These are different properties of the same ordinary map.
Their relationship is best understood by first considering full function
spaces, then restricting to supported elements. The mathematical statements
below describe the mathematics behind the function interfaces.

\subsection{Full function spaces}
\label{sec:full-function-spaces}

Let $G$ be any group acting on $X$ and $Y$. Write $[X,Y]$ for the set of
all functions $X\to Y$ with the \emph{conjugation action}
\begin{equation}
  (g\act f)(x)=g\act f(g^{-1}\act x).
  \label{eq:function-conjugation}
\end{equation}
The inverse in the argument ensures the action law:
\[
  (g\act(h\act f))(x)
    =(gh)\act f(h^{-1}\act(g^{-1}\act x))
    =((gh)\act f)(x).
\]
This is Pitts' function space for group actions
\cite[Section~1.4, equation~(1.24), p.~18]{pitts2013}.

Three operations should be distinguished. The pointwise action
$(g\act_{\mathrm{pt}}f)(x)=g\act f(x)$ renames only outputs;
inverse precomposition $(g\act_{\mathrm{pre}}f)(x)=f(g^{-1}\act x)$
renames only inputs. Conjugation combines both. It agrees with inverse
precomposition when the codomain action is trivial, and with the pointwise
action when the domain action is trivial. In general neither simplification
is valid. For example, conjugation fixes $\id_A$ on atoms, whereas a
nonidentity transposition changes it under either of the other two actions.

\begin{proposition}[Application and currying]
\label{prop:function-evaluation-curry}
For arbitrary $G$-sets $Z,X,Y$, evaluation and currying are equivariant:
\begin{gather*}
  \operatorname{ev}:[X,Y]\times X\longrightarrow Y,
  \qquad \operatorname{ev}(f,x)=f(x),\\
  \operatorname{curry}:[Z\times X,Y]\longrightarrow[Z,[X,Y]],
  \qquad \operatorname{curry}(F)(z)(x)=F(z,x).
\end{gather*}
Currying is an equivariant bijection whose inverse is uncurrying.
\end{proposition}

\begin{proof}
Evaluation satisfies
$(g\act f)(g\act x)=g\act f(x)$ by cancellation. For currying,
\[
  (g\act\operatorname{curry}(F))(z)(x)
    =g\act F(g^{-1}\act z,g^{-1}\act x)
    =\operatorname{curry}(g\act F)(z)(x).
\]
The ordinary function equations prove that currying and uncurrying are
mutual inverses; the same calculation gives equivariance of uncurrying.
\end{proof}

These calculations give the exponential universal property for group
actions \cite[Theorem~1.7, pp.~18--19]{pitts2013}. They use no atom carrier,
finite support, infinitude or nominality assumptions.

In Lean, \texttt{FunctionObject G X Y} keeps conjugation on a distinct carrier,
preserving Mathlib's pointwise action on ordinary arrows. The group parameter
selects the action interpretation without raising the universe of the function
carrier. Its ordinary-function projection supports application and extensional
equality, but is not generally equivariant to the pointwise arrow action.
This boundary permits arbitrary, possibly unsupported maps to participate in
the general action theory. Equivariant maps use Mathlib's existing
\texttt{MulActionHom} interface.

\subsection{Support of functions and their values}
\label{sec:function-support}

Now take $G=\Perm(A)$. Pitts fixes countably infinite atoms from
Chapter~2 onward; the finite-avoidance and least-support arguments here
need infinitude but not countability. For a finite atom set $S$, write
$\operatorname{SupportsMap}(S,f)$ for the ordinary-function condition
\begin{equation}
  \forall\pi\in\Perm(A),\quad
  \bigl(\forall a\in S,\ \pi(a)=a\bigr)
    \Longrightarrow
  \bigl(\forall x\in X,\ f(\pi\act x)=\pi\act f(x)\bigr).
  \label{eq:supports-map}
\end{equation}
Substituting $\pi\act x$ or $\pi^{-1}\act x$ shows that this is
exactly support of $f$ in $[X,Y]$. In particular, the empty set supports
$f$ precisely when $f$ is equivariant
\cite[Lemma~2.17, p.~35]{pitts2013}.

Let $[X,Y]_{\mathrm{fs}}$ consist of the finitely supported elements of
$[X,Y]$. Support transport makes this subset invariant under conjugation,
and the inherited action makes it a nominal set. This construction is
meaningful even when $X$ or $Y$ contains unsupported elements: it is the
supported-element construction of Pitts' Theorem~2.15, p.~34. When $X$
and $Y$ are nominal, it is the nominal function set of Definition~2.18,
p.~35 \cite{pitts2013}. Membership asserts that a finite bound exists;
it does not select an observable bound as additional function data.

The corresponding \texttt{SupportedMap A X Y} record stores support existence
as a proposition. Consequently, changing that evidence changes neither its
ordinary extensional equality nor its computational behavior. Ordinary inputs
use \texttt{SupportsMap} certificates, with a proved correspondence to support
of their full function objects. This lets a definition remain an ordinary Lean
function, while a higher-order interface can carry its support evidence with
the value. Neither interface requires every element of the domain or codomain
to be supported.

\begin{proposition}[Evaluation bound]
\label{prop:function-evaluation-support}
If $S$ supports $f\in[X,Y]$ and $T$ supports $x\in X$, then $S\cup T$
supports $f(x)$. Consequently, when $A$ is infinite and $f,x$ are
finitely supported,
\begin{equation}
  \operatorname{supp}(f(x))
    \subseteq \operatorname{supp}(f)\cup\operatorname{supp}(x).
  \label{eq:function-evaluation-bound}
\end{equation}
For a finitely supported function over infinite atoms,
$\operatorname{supp}(f)=\varnothing$ if and only if $f$ is equivariant.
\end{proposition}

\begin{proof}
A permutation fixing $S\cup T$ fixes both $f$ and $x$. Equivariance of
evaluation therefore fixes $f(x)$. Taking least supports gives the
inclusion. The final equivalence follows from
\eqref{eq:supports-map} and the characterization of empty least support.
\end{proof}

Only the particular argument $x$ needs finite support; no carrier-wide
nominality hypothesis is needed for this reasoning. In freshness form,
$a\mathbin{\#}f$ and $a\mathbin{\#}x$ imply
$a\mathbin{\#}f(x)$, as in Pitts' Proposition~3.4(ii), p.~50
\cite{pitts2013}. The converse fails: the map $T\mapsto T\setminus\{a\}$
on finite atom sets always has an output fresh for $a$, but the map itself
has least support $\{a\}$ \cite[Example~3.5, pp.~50--51]{pitts2013}.
Freshness of outputs alone cannot certify freshness of a function object.

Empty domains deserve separate attention. For a constant map $c_y:X\to Y$,
every support of $y$ supports $c_y$. If $X$ is nonempty, the converse
holds: evaluating the constant-function fixed-point equation at any
argument gives $\pi\act y=y$, without requiring that argument to be
supported. Thus supported constants have the same least support as their
values over infinite atoms. If $X$ is empty, there is only one function
$X\to Y$; it is equivariant regardless of the support of $y$, and the
constant-map construction need not be injective.

\subsection{Fixed parameters and supported currying}
\label{sec:supported-currying}

Suppose $F:Z\times X\to Y$ has finite support $S$. If $z\in Z$ has
support $T$, then the section $F_z:x\mapsto F(z,x)$ has support
$S\cup T$. Indeed, a permutation fixing this union satisfies
\[
  F(z,\pi\act x)=F(\pi\act z,\pi\act x)=\pi\act F(z,x).
\]
For jointly equivariant $F$, the parameter bound $T$ alone suffices.
Fixing an atom parameter therefore generally produces a supported function,
not an equivariant one. For example, $F(a,x)=a$ is jointly equivariant,
but for a nonempty domain its section has least support $\{a\}$ over
infinite atoms.

Full-space currying always preserves and reflects every support bound,
because it is an equivariant bijection. A different question is whether
its values belong to $[X,Y]_{\mathrm{fs}}$. For a particular $F$, the
exact condition is that every section $F_z$ is finitely supported.
Finite supportedness of $F$ together with nominality of $Z$ is a
sufficient uniform condition. Under it, currying gives a correspondence
\[
  [Z\times X,Y]_{\mathrm{fs}}
    \simeq [Z,[X,Y]_{\mathrm{fs}}]_{\mathrm{fs}}.
\]
No nominality of $X$ or $Y$ is needed for this correspondence. If all three
carriers are nominal, its restriction to equivariant maps is the
exponential correspondence in the category of nominal sets
\cite[Section~2.4 and Theorem~2.19, pp.~35--36]{pitts2013}.

The section condition cannot simply be dropped. Let $Z$ contain an
unsupported element $z$, take $X$ to be a one-point discrete set, and let
$F:Z\times X\to Z$ be the first projection. The map $F$ is equivariant,
so is finitely supported, but $F_z$ is the unsupported constant map with
value $z$. This counterexample concerns the codomain restriction of
currying; the full-space equivariant bijection remains valid.

The interface exposes this distinction through \texttt{curryAt}, using evidence
for one parameter, and \texttt{curryWithSections}, using evidence for every
section. The latter agrees with ordinary currying on application and is inverse
to uncurrying on its admissible domain. The convenient \texttt{curry} operation
derives section evidence from nominality of the parameter carrier. Support
reflection through the injective uncurrying map proves exact support equality,
rather than only an upper bound.

\subsection{Predicates with ordinary truth values}
\label{sec:function-predicates}

A predicate remains an ordinary $P:X\to\mathrm{Prop}$. Its renaming
operation is inverse precomposition:
\begin{equation}
  (\pi\act P)(x)=P(\pi^{-1}\act x).
  \label{eq:predicate-action}
\end{equation}
This is conjugation with trivial action on truth values, corresponding
to Pitts' powerset construction from a discrete two-element truth object
\cite[Section~1.5, pp.~19--20]{pitts2013}. It does not require replacing
ordinary propositions by a separate internal logic.

One can state support directly, without choosing any action instance on
\texttt{Prop}:
\begin{equation}
\begin{split}
  \operatorname{SupportsPred}(S,P)\quad\Longleftrightarrow\quad
  &\forall\pi\in\Perm(A),\quad
  \bigl(\forall a\in S,\ \pi(a)=a\bigr)\\
  &\hspace{1em}\Longrightarrow
    \bigl(\forall x,\ P(\pi\act x)\leftrightarrow P(x)\bigr).
\end{split}
\label{eq:supports-predicate}
\end{equation}
If an acted-on predicate carrier is used, function extensionality and
propositional extensionality identify this condition with support under
\eqref{eq:predicate-action}. Equivalently, closure of the satisfying set
under every permutation fixing $S$ suffices: applying the closure property
also to inverses supplies the reverse implication
\cite[Lemma~2.24, p.~39]{pitts2013}. For a supported binary predicate,
fixing a supported parameter gives the same union bound as for functions;
Pitts' Proposition~2.25, p.~39 gives this argument for subsets of products.

Support of a truth value and support of a predicate are different notions.
For example, $P_a(x)\equiv(x=a)$ is renamed to $P_{\pi a}$.
Its singleton bound $\{a\}$ supports it, but a swap moving $a$ changes it.
With trivial action each individual proposition is fixed, but its dependence
on a renamed argument can still have no finite bound. If an atom subset
$U\subseteq A$ and its complement are both infinite, the predicate
$a\mapsto a\in U$ is unsupported. Given a proposed finite support $S$,
choose $a\in U\setminus S$ and $b\notin U\cup S$. Their swap fixes $S$
and changes membership. More generally, an atom subset is finitely
supported exactly when it is finite or cofinite
\cite[Proposition~2.9, pp.~31--32]{pitts2013}. Thus arbitrary predicates
remain legitimate logical inputs, including induction motives, while their
use as supported parameters requires a separate hypothesis.
Section~\ref{sec:predicate-calculus} develops the supported predicate carrier
and its logical operations from these ordinary certificates and semantic views.

\subsection{Sufficient support from a context}
\label{sec:function-context-support}

Pitts' Equivariance Principle states that constructions from equivariant
functions and relations by higher-order logic remain equivariant
\cite[Proposition~1.9 and following principle, p.~21]{pitts2013}.
For a construction depending on parameters, the equivariance assertion
renames all those parameters together. Fixing them gives the following
useful sufficient-bound rule.

\begin{proposition}[Context support]
\label{prop:function-context-support}
Let $c_i\in C_i$ have finite support bounds $S_i$, for $1\leq i\leq n$.
If $E:(C_1\times\cdots\times C_n)\times X\to Y$ has support bound
$S_0$, then
\[
  S_0\cup S_1\cup\cdots\cup S_n
  \quad\text{supports}\quad
  \bigl(x\mapsto E((c_1,\ldots,c_n),x)\bigr).
\]
In particular, for jointly equivariant $E$ one can take $S_0=\varnothing$.
The same rule holds for predicates, interpreting the support equation
by logical equivalence.
\end{proposition}

\begin{proof}
A permutation fixing the union fixes every $c_i$. The support equation
for $E$ then specializes to
\[
  E((c_1,\ldots,c_n),\pi\act x)
    =\pi\act E((c_1,\ldots,c_n),x).
\]
For predicates the identical argument uses a biconditional instead of
equality of outputs.
\end{proof}

Only sufficient bounds for these particular values are needed. The proof
does not require least supports, infinitude of the atom carrier or
nominality of every element of each $C_i$. Consequently, explicitly
supplied support certificates suffice for this form of context reasoning;
no constructive least-support theorem is being assumed. Pitts' Finite
Support Principle gives the corresponding closure account for nominal
higher-order logic, with function and subset quantification ranging over
supported objects \cite[Section~2.5, pp.~39--40]{pitts2013}.
The distinction is consequential constructively: general existence of least
finite supports is not constructively provable, as Swan's counterexamples
show \cite[Theorem~3.5]{swan2017}. Particular structures can still have
constructively defined least supports; the context rule above only uses
the supplied sufficient bounds.

Urban's Isabelle/HOL account uses a closely related method: the free
variables of a HOL expression suggest a support bound, then swapping
equations certify it \cite[Section~5, Definition~5, Lemma~11 and
Examples~1--2]{urban2008}. His discussion of function-space integration
concerns Isabelle's infrastructure; it does not establish an obstruction
to proof-field function bundles in Lean. A suitable interface can therefore
accept ordinary functions with explicit certificates while also exposing
supported function values for higher-order nominal reasoning.

Selecting the atom carrier and its actions is distinct from establishing
these support certificates. A function type determines its domain and
codomain; it does not determine which atom parameters a particular function
uses or whether it is finitely supported. A context-based proof may combine
supplied bounds with proved equivariance rules, but merely listing local
values proves no support claim. Captured functions and predicates require
their own evidence when such a rule uses them, and fixed global values
are parameters too. This is the warning of Pitts' Note~1.10, p.~22
\cite{pitts2013}. Conversely, a bound obtained by collecting parameters
is only sufficient: an expression may ignore a parameter or cancel its
dependence, so a larger context need not describe its least support.

\subsection{The boundary of choice}
\label{sec:function-choice}

Classical choice does not remove these obligations. For infinite $A$,
there is no finitely supported function
$c:\Finset(A)\to A$ satisfying $c(S)\notin S$ for every finite $S$.
If $T$ supported such a function, then $T$ supports itself as a finite
atom set, so the evaluation bound would give
$\operatorname{supp}(c(T))\subseteq T$. The canonical atom support
$\operatorname{supp}(c(T))=\{c(T)\}$ contradicts $c(T)\notin T$.
This is the finite-avoidance form of the obstruction in Pitts'
Theorem~2.29, p.~43, concerning choice from nonempty supported atom subsets
\cite{pitts2013}. Each individual fresh atom exists, and classical Lean
may select one; the resulting global selector has no finite support.
Ordinary choice, equivariance and finite supportedness must therefore
remain separate assertions.

Pitts explicitly excludes unrestricted Hilbert choice from the
Equivariance Principle \cite[Note~1.11, p.~22]{pitts2013}. Choice used
to extract a uniquely determined value is different: if the defining
relation is equivariant and total with unique outputs, renaming preserves
the relation and uniqueness forces the chosen function's equivariance.
Thus a choice-based definition needs an appropriate mathematical argument;
the occurrence of choice alone neither supplies nor refutes its support.
